\documentclass[11pt,a4paper]{article}

%biblioteci
\usepackage[romanian]{babel}  % bibliografia si cuprinsul sa fie in romana
% \usepackage{times} % times new roman
\usepackage{csquotes} % pt warning
\usepackage{lipsum}
\usepackage{indentfirst} % pentru ca si primul paragraf sa aiba alineat
\usepackage{graphicx}  % pentru inserare imagini
\usepackage{float} % pt asezare imagini in pagina
\usepackage{hyperref} % pt cuprins cu linkuri catre sectiuni
\usepackage{minted} % pt cod
\hypersetup{
    colorlinks,
    citecolor=black,
    filecolor=black,
    linkcolor=black,
    urlcolor=black
}

\usepackage{biblatex}

% format document(margin)
\usepackage[top=2.5cm, bottom=2.5cm, left=2.5cm, right=2.5cm]{geometry}
\linespread{1.06} % spatiu intre cuvinte
\setlength{\parskip}{8pt plus2pt minus2pt} % spatiu intre paragrafe

% widow and orphan lines
\widowpenalty 10000
\clubpenalty 10000

\newcommand{\eat}[1]{}
\newcommand{\HRule}{\rule{\linewidth}{0.5mm}}

%puncte si dupa fiecare sectiune in cuprins
\usepackage{tocloft}
\renewcommand{\cftsecleader}{\cftdotfill{\cftdotsep}}

% pentru bibliografie
% \usepackage{biblatex} % Imports biblatex package
% \addbibresource{refs.bib} % Import the bibliography file

% afiseaza mai frumos listele
\usepackage{enumitem} % ceva control asupra listelor/ enumerarilor
\setlist{nolistsep,noitemsep}

\setlength{\parskip}{2pt}  % spatiu intre paragrafe

\begin{document}

\begin{titlepage}
\begin{center}

% Top 
\includegraphics[width=0.55\textwidth]{resources/utcn_logo.jpg}~\\[2cm]

% Title
\HRule \\[0.4cm]
{ \LARGE
    \textbf{Probleme de cautare si agenti adversariali}\\[0.4cm]
    \emph{Inteligenta Artificiala}\\[0.4cm]
}
\HRule \\[1.5cm]

% Author
{ \large
    Autori: Vlad Durdeu și Moroz Lion \\[0.1cm]
    Grupa: 30232\\[0.1cm]
}

\vfill
\textsc{\large Facultatea de Automatica\\si Calculatoare}\\[0.4cm]

% Bottom
{\large 2 Decembrie 2024}
    
\end{center}
\end{titlepage}
\tableofcontents % afiseaza cuprinsul
\section{Uninformed search}
\subsection{Question 1 - Depth First Search}
\par În această întrebare, scopul este găsirea unui punct fix de mâncare folosind algoritmul \textbf{Depth First Search (DFS)}. DFS explorează nodurile unui graf mergând pe cea mai adâncă cale înainte de a reveni la nodurile anterioare.

\par Codul principal este implementat în funcția \texttt{depthFirstSearch} din \texttt{search.py} și este descris mai jos:

\inputminted[linenos]{python}{code/depth_first_search.py}
\par Acest algoritm utilizează un \textit{stack} (stivă) pentru a explora nodurile și revine la cele anterioare dacă calea curentă nu este validă. Performanța sa depinde de adâncimea soluției și numărul total de noduri, având o complexitate de \( \mathcal{O}(V + E) \), unde \( V \) este numărul de noduri, iar \( E \) este numărul de muchii.


\subsection{Question 2 - Breadth First Search}
\par În această întrebare, algoritmul \textbf{Breadth First Search (BFS)} este utilizat pentru a determina calea cea mai scurtă dintre punctul de start (\texttt{problem.getStartState()}) și punctul de oprire definit de condiția \texttt{problem.isGoalState(state)}.

\par Implementarea din \texttt{search.py} este:
\inputminted[linenos]{python}{code/breadth_first_search.py}
\par  BFS explorează mai întâi nodurile la nivelul cel mai apropiat, utilizând un \textit{queue} (coadă) și având o complexitate de \( \mathcal{O}(V + E) \), unde \( V \) este numărul de noduri, iar \( E \) este numărul de muchii.

\subsection{Question 3 - Varying the Cost Function}
\par În această întrebare este prezentată funcției de cost, algoritmul fiind conceput să selecteze nodurile pe baza celui mai mic cost acumulat. 

\par Implementarea metodei este prezentată mai jos:
\inputminted[linenos]{python}{code/uniform_cost_search.py}
\par Algoritmul utilizat, \textbf{Uniform Cost Search (UCS)}, prioritizează explorarea nodurilor cu cel mai mic cost, garantând găsirea unei soluții optime în condițiile unor costuri variabile. Complexitatea algoritmului este de  \mathcal{O}(V \log V + E \log V), unde \( V \) este numărul de noduri, iar \( E \) este numărul de muchii.



\section{Informed search}
% q4 - q8 din project 1
\subsection{Question 4 - A* Search}
\par Algoritmul \textbf{A*} combină costul real \( g(n) \), care reprezintă costul de la starea inițială până la nodul curent \( n \), cu o estimare euristică \( h(n) \), ce indică costul minim estimat de la nodul curent la nodul final. Funcția de evaluare folosită este:
\[
f(n) = g(n) + h(n)
\]
unde:
\begin{itemize}
	\item \( g(n) \) este costul real acumulat până la nodul \( n \),
	\item \( h(n) \) este o estimare a costului rămas, bazată pe euristică.
\end{itemize}

\par Implementarea în cod este următoarea:
\inputminted[linenos]{python}{code/a_star_search.py}

\par Algoritmul calculează valoarea \( f(n) \) pentru fiecare nod \( n \), utilizând o coadă de priorități pentru a explora mai întâi nodurile cu cel mai mic \( f(n) \). 

\par Performanța algoritmului depinde de euristica \( h \). Dacă \( h \) este \textit{admisibilă} (nu supraestimează costul real) și \textit{consistentă} (respectă inegalitatea triunghiului), A* garantează găsirea unei soluții optime. 

\par În întrebările ce urmează, euristica utilizată este \textbf{distanța Manhattan}, definită astfel:
\[
h(n) = |x_1 - x_2| + |y_1 - y_2|
\]
unde \((x_1, y_1)\) și \((x_2, y_2)\) sunt coordonatele nodului curent, respectiv ale nodului țintă.

\subsection{Question 5 - Finding All the Corners}
\par Problema implică explorarea tuturor colțurilor dintr-un labirint. Aceasta este definită prin clasa \texttt{CornersProblem}, care include metode esențiale pentru configurarea stării inițiale, verificarea scopului și generarea succesorilor.

\par Metodele principale utilizate sunt următoarele:

\begin{itemize}
	\item \texttt{getStartState}: Această metodă definește starea inițială a problemei, incluzând poziția de start și un tuplu de patru valori booleene pentru a urmări colțurile vizitate.
	\par Starea inițială este reprezentată ca o tuplă \((\texttt{startingPosition}, \texttt{cornersVisited})\), care reprezină dacă colțurile au fost vizitate.
	\item \texttt{isGoalState}: Această metodă determină dacă starea curentă este de tip scop, adică toate cele patru colțuri au fost vizitate:
	\par Dacă toate valorile din \texttt{state[1]} sunt \texttt{True}, înseamnă că Pacman a vizitat toate colțurile, iar problema este rezolvată.
	\item \texttt{getSuccessors}: Această metodă verifică fiecare direcție posibilă (nord, sud, est, vest) și calculează poziția următoare, evitând pereții. Dacă poziția următoare este un colț, flag-ul pentru acel colț este setat la \texttt{True}.
\end{itemize}

\par Implementarea metodelor este prezentată mai jos:
\inputminted[linenos]{python}{code/finding_all_corners.py}

\subsection{Question 6 - Corners Problem: Heuristic}
\par Euristica dezvoltată pentru problema colțurilor calculează o estimare a costului minim necesar pentru a vizita toate colțurile rămase. Aceasta se bazează pe distanța Manhattan și adaugă distanțele minime necesare pentru a vizita colțurile neexplorate.

\par Implementarea metodei este următoarea:
\inputminted[linenos]{python}{code/corners_heuristic.py}

\par Această metodă funcționează astfel:
\begin{itemize}
	\item Se identifică toate colțurile care nu au fost vizitate (\texttt{unvisitedCorners}).
	\item Începând de la poziția curentă (\texttt{current}), se calculează distanța Manhattan de la nodul curent la fiecare colț neexplorat.
	\item Se selectează colțul cel mai apropiat, se adaugă distanța către acesta la euristică și se actualizează poziția curentă.
	\item Procesul continuă până când toate colțurile au fost vizitate.
\end{itemize}

\par Euristica returnează suma distanțelor minime necesare pentru a vizita toate colțurile rămase. Este:
\begin{itemize}
	\item \textbf{Admisibilă}: Nu supraestimează costul real, deoarece se folosește distanța Manhattan.
	\item \textbf{Consistentă}: Respectă inegalitatea triunghiului. Pentru orice două stări consecutive \( n \) și \( n' \), diferența dintre euristicile acestora nu depășește costul deplasării directe între ele.
\end{itemize}

\subsection{Question 7 - Eating All the Dots}
\par Scopul acestei probleme este găsirea unei soluții eficiente pentru ca Pacman să consume toate punctele de mâncare de pe hartă. Problema este definită prin clasa \texttt{FoodSearchProblem}, iar euristica utilizată este implementată în metoda \texttt{foodHeuristic}.

\par Implementarea metodei este următoarea:
\inputminted[linenos]{python}{code/eating_all_the_dots.py}

\par Această metodă funcționează astfel:
\begin{itemize}
	\item Dacă nu există puncte de mâncare (\texttt{foodList} este golă), euristica returnează 0.
	\item Se calculează următoarele valori:
	\begin{itemize}
		\item \texttt{minDistance}: Distanța Manhattan de la poziția curentă la cel mai apropiat punct de hrană.
		\item \texttt{maxDistance}: Suma diferențelor maxime între coordonatele extreme (stânga-dreapta, sus-jos) ale punctelor de mâncare.
	\end{itemize}
	\item Euristica este suma celor două valori:
	\[
	heuristic = \texttt{minDistance} + \texttt{maxDistance}
	\]
\end{itemize}

\par Această euristică are următoarele proprietăți:
\begin{itemize}
	\item Este \textbf{admisibilă}, deoarece nu supraestimează costul real. Suma distanțelor este întotdeauna mai mică sau egală cu costul minim necesar pentru a consuma toate punctele de mâncare.
	\item Este \textbf{consistentă}, deoarece respectă inegalitatea triunghiului. Pentru orice două stări consecutive \( n \) și \( n' \), diferența dintre euristicile acestora nu depășește costul deplasării directe între ele.
\end{itemize}

\subsection{Question 8 - Suboptimal Search}
\par În această problemă, se implementează o abordare suboptimă pentru consumarea tuturor punctelor de hrană. Metoda utilizată presupune identificarea celui mai apropiat punct de mâncare și navigarea către aceasta.

\par Metoda principală, \texttt{findPathToClosestDot}, este implementată astfel:
\inputminted[linenos]{python}{code/suboptimal_search.py}

\par Descrierea metodei:
\begin{itemize}
	\item Se obține poziția inițială a lui Pacman folosind \texttt{gameState.getPacmanPosition()}.
	\item Se obține grila punctelor de mâncare și grila pereților.
	\item Se creează o instanță a clasei \texttt{AnyFoodSearchProblem}, care definește problema căutării unui punct de mâncare.
	\item Se utilizează algoritmul \textbf{Breadth First Search (BFS)} pentru a găsi calea către cel mai apropiat punct de mâncare.
	\item Se mai implementează metoda \textbf{isGoalState(self, state: Tuple[int, int])} pentru a determina punctul de oprire al căutării.
\end{itemize}

\section{Adversarial search}
% q1-q3 din project2
\subsection{Question 1 - Reflex Agent}
\par Metoda \texttt{evaluationFunction} calculează un scor pentru fiecare acțiune posibilă, luând în considerare mai mulți factori ai stării jocului. Implementarea este prezentată mai jos:
\inputminted[linenos]{python}{code/reflex_agent.py}

\par Algoritmul calculează următoarele valori:
\begin{itemize}
	\item \textbf{Distanța până la cea mai apropiată mâncare} (\texttt{closestFoodDistance}): Cu cât Pacman este mai aproape de mâncare, cu atât scorul va crește.
	\item \textbf{Distanța până la capsule} (\texttt{closestCapsuleDistance}): Capsulelele oferă un scor suplimentar, încurajându-l pe Pacman să le consume.
	\item \textbf{Distanța până la fantome}: 
	\begin{itemize}
		\item Dacă o fantomă este „scared” (timp rămas pozitiv în \texttt{newScaredTimes}), apropierea de ea crește scorul.
		\item Dacă fantoma nu este „scared”, apropierea excesivă reduce scorul.
	\end{itemize}
	\item \textbf{Scorul stării succesive} (\texttt{successorGameState.getScore()}): Scorul general al stării este folosit ca bază pentru calcul.
\end{itemize}

\subsection{Question 2 - Minimax}
\par Algoritmul \textbf{Minimax} este utilizat pentru a determina cea mai bună strategie pentru Pacman într-un mediu adversarial. Pacman acționează ca \textbf{agentul max}, încercând să maximizeze scorul obținut, în timp ce fantomele, \textbf{agenții min} încearcă să reducă scorul prin minimizarea acestuia.

\par Implementarea metodei \texttt{getAction} este următoarea:
\inputminted[linenos]{python}{code/minimax.py}

\par Descrierea algoritmului:
\begin{itemize}
	\item \textbf{Punct de oprire}: Dacă starea este terminală (\texttt{win} sau \texttt{lose}) sau dacă adâncimea maximă a fost atinsă, se returnează valoarea dată de \texttt{evaluationFunction}.
	\item \textbf{Pacman (Agentul Max)}:
	\begin{itemize}
		\item Pacman caută să maximizeze scorul alegând succesorul cu cel mai mare scor minim returnat de fantome.
	\end{itemize}
	\item \textbf{Fantomele (Agenții Min)}:
	\begin{itemize}
		\item Fantomele încearcă să minimizeze scorul alegând succesorul cu cel mai mic scor returnat de Pacman.
	\end{itemize}
	\item \textbf{Trecerea între agenți}: 
	\begin{itemize}
		\item Algoritmul iterează între agenți, avansând în adâncime la fiecare tur complet.
	\end{itemize}
\end{itemize}

\subsection{Question 3 - Alpha-Beta Pruning}
\par Algoritmul \textbf{Alpha-Beta Pruning} optimizează procesul de căutare Minimax prin reducerea numărului de stări evaluate. Acesta utilizează două valori, \texttt{alpha} și \texttt{beta}, pentru a menține limitele scorurilor maxime și minime posibile în arborele de joc.

\par Implementarea metodei \texttt{getAction} este următoarea:
\inputminted[linenos]{python}{code/alpha_beta_pruning.py}

\par Descrierea algoritmului:
\begin{itemize}
	\item \textbf{Punct de oprire}: Algoritmul se oprește când jocul este câștigat, pierdut sau adâncimea maximă este atinsă, utilizând \texttt{evaluationFunction} pentru evaluare.
	\item \textbf{Alpha (\( \alpha \))}: Păstrează cel mai bun scor găsit de agentul Max (Pacman).
	\item \textbf{Beta (\( \beta \))}: Păstrează cel mai bun scor găsit de agenții Min (fantomele).
	\item \textbf{Tăierea (pruning)}: Dacă un nod are un scor mai bun pentru adversar decât limita \( \alpha \) sau \( \beta \), atunci sub-arborele acestuia este eliminat din căutare.
\end{itemize}

\subsection{Question 4 - Expectimax (BONUS)}
\par Algoritmul \textbf{Expectimax} este utilizat pentru a modela comportamentul adversarilor care acționează aleatoriu. Spre deosebire de Minimax, unde adversarii sunt asumați a fi perfect raționali, Expectimax tratează deciziile adversarilor ca alegeri probabilistice.

\par Implementarea metodei \texttt{getAction} este următoarea:
\inputminted[linenos]{python}{code/expectimax.py}

\par Descrierea algoritmului:
\begin{itemize}
    \item \textbf{Punct de oprire}: Dacă starea este terminală (\texttt{win} sau \texttt{lose}) sau dacă adâncimea maximă a fost atinsă, se returnează valoarea dată de \texttt{evaluationFunction}.
    \item \textbf{Pacman (Agentul Max)}:
    \begin{itemize}
        \item Pacman caută să maximizeze scorul pe baza succesorilor disponibili.
    \end{itemize}
    \item \textbf{Fantomele (Agenții aleatorii)}:
    \begin{itemize}
        \item Fantomele sunt modelate ca acționând uniform aleatoriu asupra acțiunilor disponibile.
        \item Scorul unui nod este calculat ca o medie ponderată a succesorilor săi.
    \end{itemize}
\end{itemize}

\par Expectimax este util pentru modelarea scenariilor unde adversarii au un comportament mai puțin predictibil, dar poate necesita resurse computaționale mai mari comparativ cu Minimax.

\subsection{Question 5 - Better Evaluation Function (BONUS)}
\par Funcția de evaluare avansată contribuie la selectarea celor mai bune acțiuni pentru Pacman pe baza unei analize a:
\begin{itemize}
    \item Distanței față de fantome (evitare sau urmărire dacă sunt „scared”).
    \item Distanței până la punctele de mâncare.
    \item Disponibilității capsulelor.
    \item Explorării zonelor nevizitate.
    \item Penalizării acțiunii \texttt{STOP}.
\end{itemize}

\par Implementarea funcției de evaluare este prezentată mai jos:
\inputminted[linenos]{python}{code/better_evaluation_function.py}

\par Descriere:
\begin{itemize}
    \item \textbf{Fantome active}: Distanța Manhattan până la fantomele apropiate este calculată, iar stările unde Pacman se află prea aproape de o fantomă activă sunt penalizate sever.
    \item \textbf{Puncte de mâncare}: Stările unde Pacman este aproape de un punct de mâncare sunt recompensate pentru a încuraja colectarea acestora.
    \item \textbf{Capsule}: Se oferă o prioritate ridicată pentru colectarea capsulelor, în special când există fantome apropiate.
    \item \textbf{Explorare}: Dacă nu există puncte de mâncare în imediata apropiere, Pacman este încurajat să exploreze zonele nevizitate.
    \item \textbf{Penalizare pentru \texttt{STOP}}: Acțiunea \texttt{STOP} este penalizată pentru a preveni inactivitatea.
\end{itemize}

\par Funcția de evaluare este esențială pentru îmbunătățirea performanței agentului, mai ales în probleme complexe cu multiple constrângeri.

 % documentatia propriu-zisa
% \printbibliography  % afisare bibliografie

\end{document} 
